\documentclass[a4paper]{article}
\usepackage{amsfonts}
\usepackage{amsmath}
\usepackage{amssymb}
\usepackage{graphicx}     

\begin{document}
  
\noindent \large Auteur: Michael Zielinski \\
\noindent \large Studierichting: Informatica\\
\noindent \large Oefeningenreeks 3F nr. 2.7\\

\medskip

\normalsize

Bepaal de schuine asymptoot van de functie: \\

$ f(x) = \dfrac{x^3 - 4x^2 - 2x + 14}{2x^2 - 4} $ \\

\bigskip

We voeren een veeltermdeling uit: \\

$ x^3 - 4x^2 - 2x + 14 = \left(2x^2 - 4\right)\left(\dfrac{x}{2} - 2\right) + 6 $ \\

Dus: \\

$ f(x) = \dfrac{\left(2x^2 - 4\right)\left(\dfrac{x}{2} - 2\right) + 6}{2x^2 - 4} $ \\

$ f(x) = \dfrac{x}{2} - 2 + \dfrac{6}{2x^2 - 4} $ \\

Omdat: \\

$ \lim\limits_{x \to \pm\infty} \dfrac{6}{2x^2 - 4} = 0 $ \\

is de schuine asymptoot: \\

$ y = \dfrac{x}{2} - 2 $ \\

\bigskip

Nu onderzoeken we de ligging ten opzichte van de asymptoot: \\

$ f(x) - A(x) = \dfrac{6}{2x^2 - 4} $ \\

$ 2x^2 - 4 = 0 $ \\

$ x^2 = 2 $ \\

$ x = \pm\sqrt{2} $ \\

\bigskip

\begin{tabular}{c|ccccc}
$x$ &  & $-\sqrt{2}$ &  & $\sqrt{2}$ &  \\ \hline
$f(x) - A(x)$ & $+$ & $|$ & $-$ & $|$ & $+$ \\
\end{tabular}

\bigskip

Voor $x \to +\infty$ geldt: \\

$ f(x) > A(x) $ \\

Voor $x \to -\infty$ geldt: \\

$ f(x) > A(x) $ \\

\begin{thebibliography}{999}
%\bibitem{a} Referentie
\end{thebibliography}
\end{document}